\section{CNAME-based tracking}
\label{sec:cname-tracking}
In this section, we provide an in-depth overview of the CNAME-based tracking ecosystem through a large-scale analysis.

\subsection{CNAME-based trackers}
An overview of the detected trackers can be found in Table~\ref{tab:trackers-list}.
For every tracker we indicated the number of publishers, counted as the number of unique eTLD+1 domains that have at least one subdomain set up to refer to a tracker (typically with a CNAME record).
Furthermore, we estimated the total number of publishers by levering DNS information from the SecurityTrails API \cite{security_trails}.
More precisely, all CNAME-based trackers either require the publishers that include them to set a CNAME record to a specific domain, or the trackers create a new subdomain for every publisher.
As such, the estimated number of publishers could be determined by finding the domains that had a CNAME record pointing to the tracker, or by listing the subdomains of the tracker domain and filtering out those that did not match the pattern that was used for publishers.
For Ingenious Technologies we were unable to estimate the total number of publishers as they use a wildcard subdomain (and thus it could not be determined whether a subdomain referred to an actual publisher using CNAME tracking).


\renewcommand{\arraystretch}{1.1}
\begin{table*}
	\caption{Overview of the analyzed CNAME-based trackers, based on the HTTP Archive dataset from October 2020.}
	\centering
	\label{tab:trackers-list}
	\begin{threeparttable}
		\adjustbox{max width=\textwidth}{
			\begin{tabular}{ M{3.8cm} r r r c c c } 
				\hline
				& & & & \multicolumn{3}{c}{\textbf{requests to tracker is blocked by}} \\
				\textbf{Tracker} & \multicolumn{1}{M{2.2cm}}{\textbf{Detected \#~publishers}} & \multicolumn{1}{M{2.2cm}}{\textbf{Est.~total \#~publishers}} & \multicolumn{1}{M{1.6cm}}{\textbf{Pricing} (min.~/mo)} & \multicolumn{1}{M{2.3cm}}{\textbf{uBlock Origin Firefox}} & \multicolumn{1}{M{2.3cm}}{\textbf{uBlock Origin Chrome}} & \multicolumn{1}{M{2.9cm}}{\textbf{NextDNS CNAME~blocklist}} \\
				\hline
				Pardot & 5,993 & 21,759 & \$1,250\hphantom{\textsuperscript{$\dagger$}} & \hphantom{\textsuperscript{*}}\faCheck\textsuperscript{*} & \hphantom{\textsuperscript{*}}\faCheck\textsuperscript{*} & \faClose \\
				Adobe Experience Cloud & 2,612 & 9,029 & \$5,000\textsuperscript{$\dagger$} & \faCheck & \faCheck & \faCheck\\ 
				Act-On Software & 1,041 & 2,533 & \$900\hphantom{\textsuperscript{$\dagger$}} & \faCheck & \faCheck & \faClose \\ 
				Oracle Eloqua & 304 & 3,743 & \$2,000\textsuperscript{$\dagger$} & \faCheck & \faClose & \faClose \\ 
				Eulerian & 253 & 1,501 & ?\hphantom{\textsuperscript{$\dagger$}} & \faCheck & \faClose & \faCheck\\
				Webtrekk & 101 & 822 & ?\hphantom{\textsuperscript{$\dagger$}} & \faCheck & \faCheck & \faCheck \\
				Ingenious Technologies & 41 & - & ?\hphantom{\textsuperscript{$\dagger$}} & \faClose & \faClose & \faCheck \\ 
				TraceDock & 49 & 69 & \euro49\hphantom{\textsuperscript{$\dagger$}} & \faClose & \faClose & \faCheck \\ 
				<intent> & 14 & 124 & ?\hphantom{\textsuperscript{$\dagger$}} & \faClose & \faClose & \faCheck \\
				AT Internet & 31 & 74 & \euro355\hphantom{\textsuperscript{$\dagger$}} & \faClose & \faClose & \faCheck\\
				Criteo & 16 & 13,082 & ?\hphantom{\textsuperscript{$\dagger$}} & \faCheck & \faClose & \faCheck \\
				Keyade & 12 & 86 & ?\hphantom{\textsuperscript{$\dagger$}} & \faCheck & \faClose & \faCheck \\
				Wizaly & 12 & 55 & \$2000\textsuperscript{$\dagger$} & \faClose & \faClose & \faCheck \\
				\hline
		\end{tabular}}
	\end{threeparttable}
	\begin{tablenotes}
		\small
		\item $\dagger$: Pricing information does not originate from original source, but as reported in reviews of the product. 
		\item *: Requests made to the CNAME subdomain triggered by a third-party analytics script hosted on \texttt{pardot.com}; the block-\\list prevents the analytics script from loading. If this script was loaded from the CNAME domain, it would not be blocked.
	\end{tablenotes}
\end{table*}

We noted the price of the services offered by the tracker suppliers when such information was available, either from the tracker's website or through third-party reviews.
In most cases, with the exception of TraceDock, which specifically focuses on providing mechanisms for circumvention of anti-tracking techniques, the offered services included a range of analytics and marketing tools.

Finally, for every tracker we determined whether tracking requests would be blocked by three relevant anti-tracking solutions: uBlock Origin (version 1.26) on both Firefox and Chrome, and the NextDNS CNAME blocklist \cite{cname_cloaking_blocklist}, which was used to extend the list of trackers we considered.
As of version 1.25 of uBlock Origin, the extension on Firefox implements a custom defense against CNAME-based tracking~\cite{ublock-origin-blocks-cname}, by resolving the domain name of requests that are originally not filtered by the standard block list and then again checks this block list against the resolved CNAME records.
Because Chrome does not support a DNS resolution API for extensions, the defense could not be applied to this browser.
Consequently, we find that four of the CNAME-based trackers (Oracle Eloqua, Eulerian, Criteo, and Keyade) are blocked by uBlock Origin on Firefox but not on the Chrome version of the anti-tracking extension.


\subsection{Tracking publishers}

\label{sec:tracking-customers}
As a result of our analysis of the HTTP Archive dataset, we detected 10,474 eTLD+1 domains that had a subdomain pointing to at least one CNAME-based tracker, with 85 publishers referring to two different trackers.
We find that for 9,501 publisher eTLD+1s the tracking request is included from a same-site origin , i.e., the publisher website has the same eTLD+1 as the subdomain it includes tracker content from.
Furthermore, on 18,451 publisher eTLD+1s we found the tracker was included from a cross-site origin; these were typically sites that were related in some way, e.g.\ belonging to the same organization.
Although these instances cannot circumvent countermeasures where all third-party cookies are blocked, e.g.\ the built-in protection of Safari, they still defeat blocklists.
\begin{figure}
	\centering
	\includegraphics[width=0.8\linewidth]{figures/stacked_tranco.pdf}
	\caption{Percentage of websites using CNAME-based tracking per bin of 10,000 ranks.}
	\label{fig:customer-ranking}
\end{figure}

Figure~\ref{fig:customer-ranking} displays the percentage of publisher eTLD+1s involved in CNAME-based tracking, both in a same-site or cross-site context, for bins of 10,000 Tranco-ranked websites.
The ratio of same-site to cross-site CNAME-based tracking is consistently between 50\% and 65\% for all bins. 
We can clearly see that the use of CNAME-based tracking is heavily biased towards more popular websites.
In the top 10,000 Tranco websites 10\% refer to a tracker via a CNAME record.
Because our dataset only contains information about the homepage of websites, and does not include results from Criteo, the reported number should be considered a lower bound.

Using the categorization service by McAfee~\cite{mcafeecategory}, we determined the most popular categories among CNAME-based tracking publishers, as shown in Figure~\ref{fig:customer-categories}.
As a baseline comparison, we also include the distribution of categories in the Tranco top 10k.
Because of the strong financial motives to perform tracking, e.g.\ marketing and attribution of online purchases, it is not surprising that publishers are mainly financially-focused, with approximately 40\% of the publisher's websites being categorized as Business.

\begin{figure}
	\centering
	\includegraphics[width=0.8\linewidth]{figures/categories_customer_websites.pdf}
	\caption{Most popular categories among CNAME-based tracking publishers.}
	\label{fig:customer-categories}
\end{figure}

Finally, we explored to what extent publishers that employ CNAME-based tracking also include third-party trackers.
To this end we analyzed all requests using the EasyPrivacy blocklist \cite{easyprivacy} to determine the number of trackers that would be blocked by this list.
We find that on the vast majority of websites that include a CNAME-based tracker (93.97\%) at least one third-party tracker was present; on average these sites had 28.43 third-party tracking requests.
This clearly shows that CNAME-based tracking is most often used in conjunction with other types of tracking.
From a privacy perspective this may cause certain issues, as the other trackers may also set first-party cookies via JavaScript; we explore this in more detail in Section~\ref{sec:implications-of-cname-tracking}.
